\begin{textsample}{NEG Dim 2 – ba – Score 1.00 – ba/ba\_ef\_6.txt}  \label{ex:f2_neg_255}
Sobre o Organizador Curricular - O DCRB elege a concepção interacionista, tratando a língua como um meio de interação social, praticada por meio dos gêneros textuais.
Essa concepção mantém estreitos laços com a concepção interacionista de aprendizagem, dando voz e vez ao sujeito em suas especificidades socioculturais, um sujeito que fala, escreve, ou vê e lê, em situações diversas e em diferentes suportes.
Nosso grande desafio é atender, de maneira significativa, a todas as modalidades de ensino, numa perspectiva inclusiva.
Cabe, então, proporcionarmos aos estudantes ( de todas as modalidades ) experiências que contribuam, em certa medida, para a ampliação dos letramentos, possibilitando -lhes a participação significativa e crítica em sociedade, utilizando sua língua. - A linguagem funciona nas diferentes práticas sociais, nas quais os sujeitos interagem por meio de textos ( orais ou escritos ) que se organizam em gêneros diversos.
Estes gêneros circulam em diferentes esferas de comunicação ou campos de atuação ( alguns em um único campo, como boletim de ocorrência ; outros circulam em vários campos, como artigo de opinião, por exemplo ).
As esferas ou campos de atuação são espaços sociais nos quais predominam determinadas atividades surgidas das necessidades humanas, como cuidar da saúde das pessoas, educar crianças, produzir literatura, divulgar acontecimentos, registrar e organizar estudos acadêmico-científicos, realizar atividades religiosas, legislar, entre outras. - Os gêneros são formas de se realizar linguisticamente objetivos específicos em situações sociais e particulares.
São, de acordo com Bakhtin²³, um conjunto de “ tipos relativamente estáveis de enunciados ”.
Marcuschi²⁴ complementa essa informação ao considerá -los como fenômenos sócio-históricos e culturalmente sensíveis ; podem ser tomados como famílias de textos com uma série de semelhança, uma vez que o predomínio da função supera a forma na determinação do gênero, evidenciando uma intertextualidade intergêneros, ou seja, um gênero com a função de outro e uma heterogeneidade tipológica quando um gênero acolhe outros vários.
Nesse sentido, o papel da escola é proporcionar aos estudantes mecanismos para que estes reconheçam as características, as finalidades e as funcionalidades destes.
Deve ainda, proporcionar aos estudantes participarem de situações de escuta e produção de textos em diferentes gêneros textuais, uma vez que esses circulam socialmente, sendo considerados como práticas sociais de diferentes esferas e campos de atividades humanas.
Outro ponto a ser considerado neste componente são os gêneros digitais, práticas de linguagem contemporâneas, que assumem, cada vez mais, características multissemióticas e multimidiáticas.
Isso se dá pelas diferentes formas de produzir, de disponibilizar, de interagir, de curtir, de \textbf{comentar} e de configurar as novas ferramentas digitais existentes e disponíveis nos ambientes da Web : redes sociais ; áudios ; vídeos ; playlists, vlogs, blogs, fanfics, podcasts ; infográficos ; enciclopédias, livros e revistas digitais, entre tantos outros.
Vale ressaltar, ainda, a importância da orientação, da checagem, da verificação dos conteúdos e publicações (re)conhecidas, na maioria das vezes e casos, como pós-verdade.
Assim, compete à escola garantir o trato, a análise, a pesquisa e a identificação das “ informações ” veiculadas na Web.
A importância da abordagem de ensino pautada nesses campos centra -se na proposta de formação integral e contextualizada com os diferentes espaços sociais : na família, na escola e nos demais espaços.
Assim sendo, os estudantes estarão inseridos em contextos de ensino e aprendizagem que os levarão à produção do conhecimento, da pesquisa e do exercício da cidadania, estabelecendo uma relação progressiva e articulada com as práticas mais cotidianas ( informais ), bem como com as mais institucionalizadas ( formais ).
As práticas de linguagens assumidas neste componente curricular dialogarão e contextualizarão com os campos de atuação previstos na BNCC.
Esses campos de atuação são utilizados para garantir que, no currículo, a escola selecione textos organizados em gêneros dos diferentes campos de atuação, em especial os gêneros de comunicação pública.
São cinco os campos de atuação considerados : campo da vida cotidiana ( somente Anos Iniciais ), campo \textbf{artístico-literário}, campo das práticas de estudo e pesquisa, campo jornalístico/midiático e campo de atuação na vida pública, sendo que esses dois últimos aparecem fundidos nos Anos Iniciais do Ensino Fundamental, com a denominação campo da vida pública : ²³ BAKHTIN, M.
Estética da criação verbal : os gêneros do discurso. 2. ed.
São Paulo : Martins Fontes, 1997. ²⁴ MARCUSCHI, L.
A.
Gêneros textuais : definição e funcionalidade.
In : DIONÍSIO, A.
Paiva et al.
Gêneros textuais \& ensino.
Rio de Janeiro : Lucerna, 2002. p. 19-36.

% matched lemmas: artístico-literário, comentar
\end{textsample}
